% End-to-end MDE pipeline of the Cyclops rover, from camera frame to range scan.
%
% Top row (left to right):    Pi frame capture -> SSH -> Depth Anything V2
%                             -> depth calibration -> back-projection to 3D.
% Bottom row (right to left): ground-plane deletion -> 1D range scan
%                             -> SSH -> Pi navigation loop.
%
% Convention (same as data-flow.tex): text above an arrow is the data carried,
% text below it is the link it travels over.
%
% The two dashed squares are image PLACEHOLDERS. To fill one, replace its
% \placeholderimg{...} with e.g. \includegraphics[width=2.2cm]{figures/xyz}.
%
% Pure vector TikZ -- solid fills and plain paths only (no opacity, shadings
% or fadings), so it also survives the latex -> dvipdf build.
\providecommand{\placeholderimg}[1]{%
  \tikz\draw[black!50, dashed] (0,0) rectangle (2.2,2.2)
    node[midway, align=center, text=black!50, font=\scriptsize] {#1};}
\begin{tikzpicture}[>=Latex, line join=round, font=\footnotesize,
    module/.style={draw=black!75, fill=white, rounded corners=2pt,
                   text width=2.2cm, align=center, minimum height=0.8cm},
    imgmodule/.style={draw=black!75, fill=white, rounded corners=2pt,
                   text width=2.4cm, align=center, inner ysep=4pt,
                   execute at begin node={\hyphenpenalty=10000}},
    imgcap/.style={font=\scriptsize\itshape, text width=2.4cm, align=center},
    group/.style={draw=black!50, fill=black!4, rounded corners=4pt},
    title/.style={font=\footnotesize\bfseries, anchor=north},
    flow/.style={->, black!80, line width=0.6pt},
    lbl/.style={align=center, inner sep=1.5pt, font=\scriptsize},
  ]

  % ---- computers (drawn first, so modules sit on top) -------------------
  \draw[group] (-0.3,-4.9) rectangle (2.3,3.6);
  \node[title] at (1.0,3.55) {Raspberry Pi};

  \draw[group] (3.7,-4.9) rectangle (14.15,3.6);
  \node[title] at (8.93,3.55) {Laptop};

  % ---- Pi modules -----------------------------------------------------------
  \node[module] (cap) at (1.0,0.6)  {Frame capture};
  \node[module] (nav) at (1.0,-3.9) {Navigation loop};

  % ---- top row: laptop --------------------------------------------------------
  \node[imgmodule] (dav2) at (5.15,0.6)
       {Depth\\Anything V2\\(\mbox{ViT-S}, frozen)\\[3pt]
        \placeholderimg{depth map\\image}\\[2pt]
        {\scriptsize\itshape Relative depth map}};

  \node[imgmodule] (cal) at (8.75,0.6)
       {Depth calibration\\[3pt]
        \placeholderimg{ground-point\\image}\\[2pt]
        {\scriptsize\itshape Ground points fit\\$a^{*}, b^{*}$ per frame}};

  \node[imgmodule] (pc) at (12.45,0.6)
       {Back-projection\\[3pt]
        \placeholderimg{point cloud\\image}\\[2pt]
        {\scriptsize\itshape 3D point cloud\\in camera frame}};

  % ---- bottom row: laptop -----------------------------------------------------
  \node[module, minimum height=1.0cm] (del) at (12.45,-3.9)
       {Ground-plane\\deletion};
  \node[module, minimum height=1.0cm] (scan) at (8.75,-3.9)
       {Reduction to\\1D range scan};

  % ---- data flow --------------------------------------------------------------
  \draw[flow] (cap.east) -- (dav2.west)
    node[lbl, midway, above] {RGB\\frame $I$}
    node[lbl, midway, below] {SSH};

  \draw[flow] (dav2.east) -- (cal.west |- dav2.east)
    node[lbl, midway, above] {$\hat{d}(u,v)$};

  \draw[flow] (cal.east) -- (pc.west)
    node[lbl, midway, above] {$\hat{z}(u,v)$};

  \draw[flow] (pc.south) -- (del.north)
    node[lbl, midway, right] {$(X,Y,Z)$};

  \draw[flow] (del.west) -- (scan.east)
    node[lbl, midway, above] {obstacle\\points};

  \draw[flow] (scan.west) -- (nav.east)
    node[lbl, pos=0.62, above] {$R_0, \dots, R_{n-1}$}
    node[lbl, pos=0.62, below] {SSH};

\end{tikzpicture}
